\documentclass[a4paper,11pt,fleqn]{article}%fleqn pour aligner les equations à gauche
\usepackage{../../Styles/Cours}


\newcommand{\chnum}{Chapitre 3}
\newcommand{\chtitle}{Cinétique chimique}
\newcommand{\dtype}{Fiche de révision}

\begin{document}
\maketitle

\section*{Ce qu'il faut connaître :}

\begin{multicols}{2}
\textbf{Le vocabulaire :}
\begin{itemize}
    \item facteurs cinétiques
    \item catalyse, catalyseur
    \item mécanise réactionnel
\end{itemize}
\columnbreak
\textbf{Les relations :}
\begin{itemize}
    \item relation entre vitesse et concentration
    \item loi de vitesse d'ordre 1
\end{itemize}
\textbf{Les grandeurs :}
\begin{itemize}
    \item vitesse volumique de disparition d'un réactif ou d'apparition d'un produit
    \item temps de demi-réaction
\end{itemize}
\end{multicols}

\section*{Ce qu'il faut savoir faire :}
\begin{tabular}{|m{.8\linewidth}|>{\centering\arraybackslash}m{.15\linewidth}|}
    \hline \centering \textbf{Capacité} & \textbf{Où la trouver}\\
    \hline Justifier le choix d’un capteur de suivi temporel de l’évolution d’un système. & \\
    \hline Identifier, à partir de données expérimentales, des facteurs cinétiques. & \\
    \hline Citer les propriétés d’un catalyseur et identifier un catalyseur à partir de données expérimentales. & \\
    \hline Mettre en évidence des facteurs cinétiques et l’effet d’un catalyseur.& \\
    \hline À partir de données expérimentales, déterminer une vitesse volumique de disparition d’un réactif, une vitesse volumique d’apparition d’un produit ou un temps de demi-réaction. & \\
    \hline Mettre en œuvre une méthode physique pour suivre l’évolution d’une concentration et déterminer la vitesse volumique de formation d’un produit ou de disparition d’un réactif. & \\
    \hline Identifier, à partir de données expérimentales, si l’évolution d’une concentration suit ou non une loi de vitesse d’ordre 1. & \\
    \hline \underline{Capacité numérique :} À l’aide d’un langage de programmation et à partir de données expérimentales, tracer l’évolution temporelle d’une concentration, d’une vitesse volumique d’apparition ou de disparition et tester une relation donnée entre la vitesse volumique de disparition et
la concentration d’un réactif. & \\
    \hline A partir d'un mécanisme réactionnel fourni, identifier un intermédiaire réactionnel, un catalyseur et établir l'équation de la réaction qu'il modélise au niveau microscopique & \\
    \hline Représenterles flèches courbes d'un acte élémentaire, en justifiant leur sens. & \\
    \hline Interpréter l'influence des concentrations et de la température sur la vitesse d'un acte élémentaire, en termes de fréquence et d'efficacité des chocs entre entités&\\
    \hline
    \end{tabular}
\end{document}